% 11-711 Fall 2026, Homework 2 handout.
% Typeset from HANDOUT.md; keep the two in sync when either changes.

\documentclass[11pt]{article}

\usepackage[margin=1in]{geometry}
\usepackage{xcolor}
\usepackage{titlesec}
\usepackage{tcolorbox}
\usepackage{amsmath,amssymb}
\usepackage{hyperref}
\usepackage{enumitem}
\usepackage{fancyhdr}
\usepackage{lmodern}

\tcbuselibrary{skins}

\definecolor{mauve}{rgb}{0.58,0,0.82}
\definecolor{headerblue}{rgb}{0.85,0.87,0.98}
\definecolor{deliverblue}{rgb}{0.06,0.10,0.45}
\definecolor{lightbg}{rgb}{0.96,0.97,1.0}
\definecolor{tabgold}{rgb}{0.99,0.95,0.80}
\definecolor{codebg}{rgb}{0.97,0.97,0.97}

\hypersetup{colorlinks=true, linkcolor=blue, urlcolor=blue, citecolor=blue}

% "Section: Title" banner
\newcommand{\problem}[1]{%
  \vspace{1em}
  \noindent\colorbox{headerblue}{\parbox{\dimexpr\textwidth-2\fboxsep}{%
    \vspace{2pt}\textbf{\large #1}\vspace{2pt}}}
  \vspace{0.6em}
}

% Subsection heading inside a banner section
\newcommand{\subsec}[1]{%
  \vspace{0.8em}
  \noindent{\sffamily\bfseries\color{mauve} #1}
  \vspace{0.35em}
  \par
}

% Boxed deliverables callout with a gold tab
\newtcolorbox{deliverables}[1]{
  enhanced, colback=lightbg, colframe=deliverblue, boxrule=0.9pt,
  arc=2pt, left=8pt, right=8pt, top=12pt, bottom=8pt,
  overlay={\node[fill=tabgold, draw=deliverblue, line width=0.9pt,
    rounded corners=1pt, inner xsep=8pt, inner ysep=3pt,
    font=\bfseries\footnotesize, anchor=north]
    at ([yshift=6pt]frame.north) {#1};},
  before skip=14pt, after skip=10pt
}

% Verbatim block on a light ground
\newtcolorbox{codebox}{
  colback=codebg, colframe=codebg!70!black, boxrule=0.4pt, arc=1pt,
  left=6pt, right=6pt, top=4pt, bottom=4pt,
  before skip=8pt, after skip=8pt
}

\pagestyle{fancy}
\fancyhf{}
\lhead{\footnotesize 11-711 Fall 2026; Homework 2}
\rhead{\footnotesize\thepage}
\renewcommand{\headrulewidth}{0.4pt}

% Long \texttt{} URLs cannot hyphenate; let TeX stretch interword space rather
% than overflow the right margin.
\emergencystretch=3em
\sloppy
\hbadness=10000

\begin{document}

\begin{center}
  {\Huge\sffamily\color{mauve} Homework 2}\\[6pt]
  {\Large\sffamily End-to-end NLP System (RAG) Building}\\[10pt]
  {\large 11-711 Advanced Natural Language Processing, Fall 2026}\\[4pt]
  \textbf{Released:} Friday, September 25 \quad\textbf{Due:} Friday, October 23, 11:59 PM ET\\[2pt]
  \textbf{Individual assignment} \quad 100 points \quad 10\% of your grade
\end{center}

\vspace{0.8em}
\hrule
\vspace{0.8em}

\noindent Large language models (LLMs) such as Llama\,2 have been shown effective for
question-answering \mbox{(\href{https://arxiv.org/abs/2307.09288}{Touvron et al., 2023})}, however,
they are often limited by their knowledge in certain domains. A common technique here is to augment
an LLM's knowledge with documents that are relevant to the question. In this assignment, you will
\emph{develop a retrieval augmented generation system (RAG)}
\mbox{(\href{https://arxiv.org/abs/2005.11401}{Lewis et al., 2021})} that is capable of answering
questions about Pittsburgh and CMU, including history, culture, trivia, and upcoming
events.\footnote{See the \href{https://cmu-anlp.org/assignments/}{assignment policies} for this
class, including submission information, late day policy and more.}

\begin{codebox}
\begin{verbatim}
Q: Who is Pittsburgh named after?
A: William Pitt

Q: What famous machine learning venue had its first conference
   in Pittsburgh in 1980?
A: ICML

Q: What musical artist is performing at PPG Arena on October 13?
A: Billie Eilish
\end{verbatim}
\end{codebox}

\noindent So far in your machine learning classes, you may have experimented with standardized tasks
and datasets that were easily accessible. However, in the real world, NLP practitioners often have to
solve a problem from scratch (like this one!). This includes gathering and cleaning data, annotating
your data, choosing a model, iterating on the model, and possibly going back to change your data. In
this assignment, you'll get to experience this full process.

\vspace{0.6em}
\noindent\textcolor{red}{\textbf{There is no starter code.}} You must collect your own data and
develop a model of your choice on the data. We provide a public leaderboard for you to evaluate your
system before the final unseen test set. We will be releasing the inputs for the test set a few days
before the assignment deadline, and you will run your already-constructed system over this data and
submit the results. We also ask you to follow several experimental best practices, and describe the
result in your report.

\vspace{0.6em}
\noindent\textbf{The key checkpoints for this assignment are:}
\begin{itemize}[nosep, leftmargin=2em]
  \item Understand the task specification
  \item Prepare your raw data
  \item Develop a retrieval augmented generation system with retrieval components
  \item Generate results
  \item Write a report
  \item Submit your work
\end{itemize}

\vspace{0.6em}
\noindent\textbf{On the use of AI.} Per the syllabus, you may use AI to help you write \emph{code}.
You may \textbf{not} use it to write any part of your report.

%======================================================================
\problem{Task: Retrieval Augmented Generation (RAG)}

\noindent You'll be working on the task of factual question-answering (QA). We will focus
specifically on questions about various facts concerning Pittsburgh and CMU. Since existing QA
systems might not have the necessary knowledge in this domain, you will need to augment each question
with relevant documents. Given an input question, your system will first retrieve documents and use
those documents to generate an answer.

%======================================================================
\problem{Preparing raw data}

\subsec{Compiling a knowledge resource}

\noindent For your test set and the RAG systems, you will first need to compile a knowledge resource
of relevant documents. We have provided a \textbf{baseline knowledge resource} that you may use as a
starting point. You can build on top of this baseline by adding new sources and expanding coverage.
You are free to use any publicly available resource, but we highly recommend including the websites
below. Note that we can also ask you questions from relevant subpages (e.g. ``about'', ``schedule'',
``history'', ``upcoming events'', ``vendors'', etc.) from these websites:

\vspace{0.4em}
\noindent\textbf{General Info and History of Pittsburgh/CMU}
\begin{itemize}[nosep, leftmargin=2em]
  \item Wikipedia pages (\href{https://en.wikipedia.org/wiki/Pittsburgh}{Pittsburgh},
    \href{https://en.wikipedia.org/wiki/History_of_Pittsburgh}{History of Pittsburgh}).
  \item \href{https://www.pittsburghpa.gov/Home}{City of Pittsburgh webpage}
  \item \href{https://www.britannica.com/place/Pittsburgh}{Encyclopedia Brittanica page}
  \item \href{https://www.visitpittsburgh.com}{Visit Pittsburgh webpage}: This website also contains
    subpages that would be useful for other topics (see below), like events, sports, music, food, etc.
  \item City of Pittsburgh \href{https://pittsburghpa.gov/finance/tax-forms}{Tax Regulations}: See the
    links under the ``Regulations'' column of the table
  \item City of Pittsburgh
    \href{https://www.pittsburghpa.gov/files/assets/city/v/4/omb/documents/operating-budgets/2025-operating-budget.pdf}{2025
    Operating Budget}
  \item \href{https://www.cmu.edu/about/}{About CMU \& CMU History}
\end{itemize}

\vspace{0.4em}
\noindent\textbf{Events in Pittsburgh and CMU} (We will only ask about annual/recurring events and
events taking place from \textbf{October 2026} onward.)
\begin{itemize}[nosep, leftmargin=2em]
  \item \href{https://pittsburgh.events}{Pittsburgh events calendar}: Navigate to month-specific pages
    for easier scraping
  \item \href{https://downtownpittsburgh.com/events/}{Downtown Pittsburgh events calendar}
  \item \href{https://www.pghcitypaper.com/pittsburgh/EventSearch?v=d}{Pittsburgh City Paper events}
  \item \href{https://events.cmu.edu}{CMU events calendar} and
    \href{https://www.cmu.edu/engage/alumni/events/campus/index.html}{campus events page}
\end{itemize}

\vspace{0.4em}
\noindent\textbf{Music and Culture} (Note that many of these pages also contain upcoming events, also
see Wikipedia pages for each.)
\begin{itemize}[nosep, leftmargin=2em]
  \item Pittsburgh \href{https://www.pittsburghsymphony.org}{Symphony},
    \href{https://pittsburghopera.org}{Opera}, and \href{https://trustarts.org}{Cultural Trust}
  \item Pittsburgh Museums (\href{https://carnegiemuseums.org}{Carnegie Museums},
    \href{https://www.heinzhistorycenter.org}{Heinz History Center}),
    \href{https://www.thefrickpittsburgh.org}{The Frick}, and
    \href{https://en.wikipedia.org/wiki/List_of_museums_in_Pittsburgh}{more}
\end{itemize}

\vspace{0.4em}
\noindent\textbf{Food-related events}
\begin{itemize}[nosep, leftmargin=2em]
  \item \href{https://www.visitpittsburgh.com/events-festivals/food-festivals/}{Food Festivals}
  \item \href{https://www.picklesburgh.com/}{Picklesburgh}
  \item \href{https://www.pghtacofest.com/}{Pittsburgh Taco Fest}
  \item \href{https://pittsburghrestaurantweek.com/}{Pittsburgh Restaurant Week}
  \item \href{https://littleitalydays.com}{Little Italy Days}
  \item \href{https://bananasplitfest.com}{Banana Split Fest}
\end{itemize}

\vspace{0.4em}
\noindent\textbf{Sports} (Note that many of these pages also contain upcoming events, also see
Wikipedia pages for each. Don't worry about scraping news/scores/recent stats from these sites.)
\begin{itemize}[nosep, leftmargin=2em]
  \item General info
    (\href{https://www.visitpittsburgh.com/things-to-do/pittsburgh-sports-teams/}{Visit Pittsburgh})
  \item Pittsburgh \href{https://www.mlb.com/pirates}{Pirates},
    \href{https://www.steelers.com}{Steelers}, and \href{https://www.nhl.com/penguins/}{Penguins}
\end{itemize}

\subsec{Collecting raw data}

\noindent Your knowledge resource might include a mix of HTML pages, PDFs, and plain text documents.
You will need to clean this data and convert it into a file format that facilitates your model
development. Here are some tools that you could use:

\begin{itemize}[nosep, leftmargin=2em]
  \item To process HTML pages, you can use
    \href{https://pypi.org/project/beautifulsoup4/}{beautifulsoup4}.
  \item To parse PDF documents into plain text, you can use
    \href{https://github.com/py-pdf/pypdf}{pypdf} or
    \href{https://github.com/jsvine/pdfplumber}{pdfplumber}.
\end{itemize}

\noindent By the end of this step, you will have a collection of documents that will serve as the
knowledge resource for your RAG system.

%======================================================================
\problem{Training data}

\noindent This is an \textbf{optional} aspect of this assignment, should you choose to train a model.
The choice of training data is a bit more flexible, and depends on your implementation. If you are
fine-tuning a model, you could possibly:

\begin{itemize}[nosep, leftmargin=2em]
  \item Annotate it yourself manually, see the guidelines below.
  \item Do some sort of automatic annotation and/or data augmentation. You cannot use closed-source
    models (such as GPT\,5, Claude etc.) for this process.
  \item Use existing datasets for transfer learning.
\end{itemize}

\noindent If you are using an LLM in a few-shot learning setting, you could possibly:

\begin{itemize}[nosep, leftmargin=2em]
  \item Annotate examples for the task using the guidelines below.
  \item Use existing datasets to identify examples for in-context learning.
\end{itemize}

\noindent This training set you have constructed will constitute
\texttt{data/train/questions.txt} and \texttt{data/train/reference\_answers.json} in your submission.
Read our model and data policy for this assignment.

\subsec{Annotation guidelines (optional)}

\begin{itemize}[nosep, leftmargin=2em]
  \item \textbf{Domain Relevance}: Training examples should be similar in nature to the evaluation
    data (i.e., questions about Pittsburgh and CMU). Use the knowledge resources listed above when
    curating examples.
  \item \textbf{Diversity}: Your training data should cover a wide range of questions about
    Pittsburgh and CMU.
  \item \textbf{Size}: The size of your training data should be sufficient to support your training
    objective (e.g., fine-tuning or prompt selection). There is no minimum required size.
  \item If you want some guidelines about this, see the lecture on experimental design and
    evaluation.\footnote{Two lectures this semester cover this material: \textbf{Research Skills and
    Experimental Design} (October 8) and \textbf{Benchmarking and Evaluation Techniques} (October 22).
    Slides are posted on the \href{https://cmu-anlp.org/schedule/}{course schedule}. Since the second
    falls the day before this assignment is due, you may also find the Fall 2025 notes on
    \href{https://cmu-l3.github.io/anlp-fall2025/static_files/anlp-f2025-14-experimentation.pdf}{experimental
    design} and
    \href{https://cmu-l3.github.io/anlp-fall2025/static_files/anlp-f2025-13-evaluation.pdf}{evaluation}
    useful earlier in the assignment.}
  \item \textbf{Quality}: Annotated examples should be accurate.
\end{itemize}

\noindent To help you get started, here are some example questions:

\begin{itemize}[nosep, leftmargin=2em]
  \item Questions that could be answered by just prompting an LLM
    \begin{itemize}[nosep, leftmargin=1.6em]
      \item Which three rivers meet in downtown Pittsburgh?
    \end{itemize}
  \item Questions that can be better answered by augmenting an LLM with relevant documents
    \begin{itemize}[nosep, leftmargin=1.6em]
      \item What is the name of the annual pickle festival held in Pittsburgh?
    \end{itemize}
  \item Questions that are likely answered only through augmentation
    \begin{itemize}[nosep, leftmargin=1.6em]
      \item When was the Pittsburgh Soul Food Festival established?
    \end{itemize}
  \item Questions that are sensitive to temporal signals
    \begin{itemize}[nosep, leftmargin=1.6em]
      \item Who is performing at X venue on Y date?
    \end{itemize}
\end{itemize}

%======================================================================
\problem{Developing your RAG system}

\noindent Unlike assignment 1, there is no starter code for this assignment. You are free to use open
source libraries and models. However, keep in mind the specific constraints outlined across different
sections in this assignment, and make sure you provide due credit for any resources used in your
report. See our model policy.

\subsec{Implementation Requirements}

\noindent For your RAG system, you will need the following three components:

\begin{enumerate}[nosep, leftmargin=2em]
  \item \textbf{Document \& query embedder} (can use existing models)
  \item \textbf{Document retriever} (implement sparse, dense and hybrid retrieval)
  \item \textbf{Document reader} (aka. question-answering system) (can use existing models)
\end{enumerate}

\vspace{0.5em}
\noindent For the core retrieval components listed below, you \textbf{MUST} implement the retrieval
logic directly using libraries such as sentence-transformers, scikit-learn, numpy etc. You cannot use
high-level RAG frameworks like LangChain, LlamaIndex, or similar tools that abstract away the
retrieval implementation.

\vspace{0.6em}
\noindent\textbf{Document Chunking.} Large documents need to be broken down into smaller, manageable
pieces (chunks) for effective retrieval, since embedding models have token limits and retrieving
overly large text segments can dilute relevance and overwhelm the generation model. Implement
document chunking strategies that split documents into appropriately-sized segments while preserving
semantic coherence. You may want to implement and try out different approaches (e.g., fixed-size with
overlap, sentence-aware chunking, paragraph-based chunking etc.), since chunking strategy can have a
significant impact on downstream retrieval and generation performance. Consider edge cases like very
short documents, very long paragraphs, and documents with special formatting.

\vspace{0.6em}
\noindent\textbf{Hybrid Retrieval.} Implement a system that combines dense (vector-based) and sparse
(keyword-based) retrieval:

\begin{itemize}[nosep, leftmargin=2em]
  \item \textbf{Dense Retrieval}: Use embedding models to create vector representations of documents
    and queries. We recommend starting with sentence-transformers models like
    \texttt{all-MiniLM-L6-v2} for efficiency, though you can experiment with other models (for
    guidance on selecting high-performing embedding models, refer to the
    \href{https://huggingface.co/spaces/mteb/leaderboard}{MTEB Leaderboard}) to improve performance.
    Use \href{https://github.com/facebookresearch/faiss}{FAISS} for efficient vector storage and
    similarity search. You'll need to implement the pipeline for embedding documents, building the
    FAISS index, and retrieving similar documents for queries.
  \item \textbf{Sparse Retrieval}: Integrate BM25 into your pipeline. You can use existing BM25
    implementations, limited to \href{https://github.com/xhluca/bm25s}{bm25s} or
    \href{https://pypi.org/project/rank-bm25/}{rank-bm25}, or build your own implementation if you
    wish to do so.
  \item \textbf{Combination Strategy}: Implement at least one established method for combining dense
    and sparse retrieval results. Common approaches include score normalization and weighted
    averaging, rank-based fusion methods (like Reciprocal Rank Fusion), using one method to filter or
    re-rank results from the other, or other ensemble techniques found in the literature. Experiment
    with different combination strategies to see which works best for your dataset.
\end{itemize}

\vspace{0.5em}
\noindent To get started with the overall RAG architecture, you can reference:

\begin{itemize}[nosep, leftmargin=2em]
  \item 11-711 lecture notes
  \item \href{https://acl2023-retrieval-lm.github.io}{ACL 2023 tutorial on retrieval-augmented LMs}
  \item \href{https://github.com/facebookresearch/llama-recipes/tree/main/demo_apps/RAG_Chatbot_example}{llama-recipes}
    for an example RAG chatbot with Llama\,2
  \item \href{https://github.com/ollama/ollama}{Ollama} or
    \href{https://github.com/ggerganov/llama.cpp}{llama.cpp} to run LLMs locally on your machine
\end{itemize}

\noindent All the code for your data preprocessing, model development and evaluation will be a part of
your GitHub repository (see submission for details).

%======================================================================
\problem{Generating results}

\noindent Finally, you will evaluate your systems on held-out test queries and submit your results for
scoring. To support iterative development, we provide a public leaderboard to help you improve your
system iteratively, as well as a separate unseen test set that will be used for final grading.

\vspace{0.5em}
\noindent The final unseen test set (questions only) will be released the day before the assignment is
due (\textbf{Thursday, October 22nd, 2026}).

\subsec{Public evaluation data}

\noindent We release \texttt{public\_evaluation\_data.json} alongside this handout: 157 questions
about Pittsburgh and CMU, each with one or more reference answers, drawn from the same pool as the
test set. Each entry looks like this:

\begin{codebox}
\begin{verbatim}
{
    "id": "1",
    "question": "What open mic event occurs weekly on Mondays
                 from 8:00-10:00 PM?",
    "answers": ["Monday Night Open Mic Comedy!"]
}
\end{verbatim}
\end{codebox}

\noindent Because the reference answers are included, you can score this set locally with the metrics
described below. Use it as your development set: it is how you are expected to run the comparisons
that the analysis section asks for. Do not train or tune on it if you also plan to report it as a
held-out number.

\vspace{0.5em}
\noindent Note that this is \textbf{not} the leaderboard set. The leaderboard uses a different set of
queries, and we will not release its reference answers.

\subsec{Public leaderboard}

\noindent\textcolor{red}{\textbf{The leaderboard is not live yet --- we will post the leaderboard
link along with queries (without answers) on Piazza as soon as it is released.}}

\vspace{0.4em}
\noindent For this leaderboard:
\begin{itemize}[nosep, leftmargin=2em]
  \item The queries are \textbf{different from} those in \texttt{public\_evaluation\_data.json}.
    Scoring well on the public evaluation data does not guarantee the same score here.
  \item We will \textbf{not} release the reference answers for the leaderboard queries, so you cannot
    score this set yourself.
  \item Only aggregate scores will be shown.
  \item You are limited to \textbf{at most 10 submissions} in total.
  \item You are only allowed to submit \textbf{one answer per question}.
  \item If you don't want your Andrew ID displayed, you can replace it with the nickname we assigned
    to you through email titled ``Advanced NLP Assignment 2 Leaderboard Nicknames''.
\end{itemize}

\vspace{0.4em}
\noindent The leaderboard is intended solely as a development and diagnostic tool. Scores on the
leaderboard will \textbf{not} be used directly for grading.

\begin{deliverables}{SUBMISSION FORMAT: PUBLIC LEADERBOARD}
\begin{verbatim}
{
    "andrewid": <your Andrew ID or assigned nickname>,
    "1": "Answer 1",
    "2": "Answer 2",
    ...
}
\end{verbatim}
\textbf{Please make sure you follow this format. Points will be deducted for misformatted outputs.}
\end{deliverables}

\subsec{Running your ablations}

\noindent The leaderboard is \textbf{not} the place to run your ablations: you have at most 10
submissions in total, and the analysis section asks you to compare several systems (dense-only,
sparse-only, hybrid, more than one fusion strategy, and retrieval vs closed-book). That is already
more configurations than you have submissions.

\vspace{0.5em}
\noindent Instead, run every comparison locally against \texttt{public\_evaluation\_data.json},
which ships with reference answers for exactly this purpose. If you annotate your own questions,
holding out a portion of them gives you a second development set; never train, tune, or prompt-select
on whatever you report as held-out. Report how many questions each number is computed over. Your
development scores and the leaderboard scores measure different question sets, so keep them in
separate columns of your results table and never average the two.

\vspace{0.5em}
\noindent Use your leaderboard submissions sparingly. Once your ablations are done, submit
predictions from the systems that scored best on the public evaluation data and on any test data you
built yourself, to check them against questions you did not write.

\subsec{Unseen test set}

\noindent This test set will be curated by the course staff and will serve as the \textbf{final
evaluation} of your system's ability to respond to a variety of questions about Pittsburgh and CMU.
Because the goal of this assignment is not to perform hyperparameter optimization on this private test
set, we ask you to not overfit to this test set. You are allowed to submit up to three output files
(\texttt{system\_outputs/system\_output\_\{1,2,3\}.json}). We will use the best performing file for
grading.

\begin{deliverables}{SUBMISSION FORMAT: UNSEEN TEST SET}
\begin{verbatim}
{
    "1": "Answer 1",
    "2": "Answer 2",
    ...
}
\end{verbatim}
\textbf{Please make sure you follow this format. Points will be deducted for misformatted outputs.}
\end{deliverables}

\subsec{Release policy for late days}

\noindent The \href{https://cmu-anlp.org/syllabus/}{course late policy} applies: for each
unexcused day your homework is late, 5\% is subtracted from your grade for this homework. You may
request excused late days by emailing \texttt{f2026-11-711@andrew.cmu.edu} at least 168 hours
(7 days) before the deadline; medical requests with documentation can be made later.

\vspace{0.5em}
\noindent This assignment adds one logistical requirement on top of that policy. Because the test set
is released the day before the deadline, everyone who submits late must receive a \textbf{different}
test set, so that students using late days do not gain extra time on the same questions. If you plan
to submit late, you \textbf{MUST} tell the TAs via a Piazza post how many days you plan to take and
the date you plan to submit. A day before that date, a different test set will be released to you. We
have distinct test sets prepared for up to 5 late days.

\subsec{Evaluation metrics}

\noindent Your submissions will be evaluated on a combination of metrics, including standard metrics
(answer recall, F1, and ROUGE-L) and LLM-as-Judge (score 1--5) with rubrics. See section 6.1 of the
original SQuAD paper for details of standard metrics. The standard metrics are token-based and measure
the overlap between your system answer and the reference answer(s). Therefore, we recommend keeping
your system generated responses as concise as possible.

%======================================================================
\problem{Writing your report}

\noindent We ask you to write a report detailing various aspects about your end-to-end system
development (see the grading criteria below).

\vspace{0.4em}
\noindent There will be a \textbf{7 page limit} for the report, and we require you to use the
\href{https://github.com/acl-org/acl-style-files}{ACL template}.

\vspace{0.4em}
\noindent Make sure you cite all your sources (open-source models, libraries, papers, blogs etc.) in
your report.

%======================================================================
\problem{Submission \& Grading}

\begin{deliverables}{SUBMISSION CHECKLIST}
There are \textbf{two Gradescope submissions}, both due at the deadline above:
\begin{itemize}[nosep, leftmargin=1.6em]
  \item \textbf{Assignment2: Report} --- \texttt{report.pdf}, built from the ACL template.
  \item \textbf{Assignment2: Code} --- a zip containing your code link, system outputs, and
    (optionally) your annotated training data.
\end{itemize}
\vspace{0.4em}
Between them you must turn in:
\begin{itemize}[nosep, leftmargin=1.6em]
  \item Your report.
  \item A link to your GitHub repository containing your code.\footnote{Create a private GitHub repo
    and give access to the TAs in charge of this assignment by the deadline. See the Piazza
    announcement post for our GitHub usernames.}
  \item (Optionally) training data you annotated for this assignment.
  \item Your system outputs on our test set.
\end{itemize}
\end{deliverables}

\noindent The zip you upload to \textbf{Assignment2: Code} should have the following structure
(assuming the lowercase Andrew ID is \texttt{ANDREWID}).

\begin{codebox}
\begin{verbatim}
ANDREWID/
|-- github.txt
|-- data/
|   |-- train (optional)/
|       |-- questions.txt
|       |-- reference_answers.json
|-- system_outputs/
|   |-- system_output_1.json
|   |-- system_output_2.json (optional)
|   |-- system_output_3.json (optional)
|-- README.md
\end{verbatim}
\end{codebox}

\noindent Do not put \texttt{report.pdf} inside the zip; it is graded on its own Gradescope
assignment.

\vspace{0.5em}
\noindent\textcolor{red}{\textbf{A code submission without \texttt{system\_outputs/} will be marked
incomplete.}} Your system outputs are what the 30 results points are scored on, so a zip that
contains only code and a GitHub link cannot be graded.

\subsec{Grading}

\noindent The following points (max.\ 100 points) are derived from the results and your report. See
the course grading policy.\footnote{Grading policy: \url{https://cmu-anlp.org/syllabus/}}

\begin{itemize}[leftmargin=2em]
  \item \textbf{Submit code (30 points)}: submit your code for preprocessing and model development in
    the form of a GitHub repo. We may not necessarily run your code, but we will look at it. So please
    ensure that it contains up-to-date code with a README file outlining the steps to run it.
  \item \textbf{Results (30 points)}: points based on your system's performance on our private test
    set. 20 points based on your performance using our metrics,\footnote{In general, if your system is
    generating answers that are relevant to the question, it would be considered non-trivial. This
    could be achieved with a basic RAG system.} plus up to 10 points based on level of performance
    relative to other submissions from the class.
  \item \textbf{Report}: the points below are awarded based on your report.
    \begin{itemize}[leftmargin=1.6em]
      \item \textbf{Data creation (10 points)}: clearly describe how you created your data. Please
        include the following details:
        \begin{itemize}[nosep, leftmargin=1.4em]
          \item How did you compile your knowledge resource, and how did you decide which documents to
            include?
          \item How did you extract raw data? What tools did you use?
          \item What data was annotated for training if any (what kind and how much)?
          \item For training data that you did not annotate, did you use any extra data and in what
            way?
        \end{itemize}
      \item \textbf{Model details (10 points)}: clearly describe your model(s). Please include the
        following details:
        \begin{itemize}[nosep, leftmargin=1.4em]
          \item What kind of methods (including baselines) did you try? Explain at least two
            variations (more is welcome). This can include variations of models, which data it was
            trained on, training strategy, embedding models, retrievers, re-rankers, etc.
          \item What was your justification for trying these methods?
        \end{itemize}
      \item \textbf{Results (10 points)}: report raw numbers from your experiments. Please include the
        following details:
        \begin{itemize}[nosep, leftmargin=1.4em]
          \item What was the result of each system variant you tried, measured on
            \texttt{public\_evaluation\_data.json} and on any development or test data you built
            yourself? Say which set each number comes from and how many questions it covers.
          \item For the systems whose predictions you submitted to the leaderboard, what score did
            they receive? You have at most 10 submissions, so we do not expect a leaderboard score
            for every variant.
        \end{itemize}
      \item \textbf{Analysis (10 points)}: perform quantitative/qualitative analysis and present your
        findings:
        \begin{itemize}[nosep, leftmargin=1.4em]
          \item Perform a comparison of the outputs on a more fine-grained level than just holistic
            accuracy numbers, and report the results. For instance, how did your models perform across
            various types of questions?
          \item Report your results across at least two variations you tried, including variations of
            models, which data it was trained on, training strategy, embedding models, retrievers,
            re-rankers, etc.
          \item Perform an analysis that evaluates the effectiveness of the retrieve-and-augment
            strategy vs closed-book use of your models.
          \item Evaluate your hybrid retrieval approach by comparing dense-only, sparse-only, and
            hybrid retrieval performance. Which fusion strategies work best for different types of
            questions?
          \item Show examples of outputs from at least two of the systems you created. Ideally, these
            examples could be representative of the quantitative differences that you found above.
        \end{itemize}
    \end{itemize}
\end{itemize}

%======================================================================
\problem{Model and Data Policy}

\noindent To make the assignment accessible to everyone:

\begin{itemize}[leftmargin=2em]
  \item You are only allowed to use models that are also accessible through
    \href{https://huggingface.co/models}{HuggingFace}. This means you may not use closed models like
    OpenAI models, but you can opt to use a hosting service for an open model (such as the Hugging
    Face or Together APIs). \textbf{The model must have been released before January 1, 2026, and must
    not exceed 32B total parameters}. For mixture-of-experts models this is the total parameter
    count across all experts, not the active count per token.
  \item You are only allowed to include publicly available data in your knowledge resource and
    training data.
  \item You are welcome to use any open-source library to assist your data annotation and model
    training. For data annotation, you can use tools like Label Studio, Doccano, or similar annotation
    platforms to create your question-answer pairs efficiently. For model development, you can use
    standard ML libraries like scikit-learn, PyTorch, or HuggingFace Transformers for any model
    training, fine-tuning, or evaluation tasks. Make sure you check the license and provide due credit
    for all tools used.
\end{itemize}

\noindent If you have any questions about whether a model or data is allowed, please ask on Piazza.

%======================================================================
\problem{FAQ}

\noindent\textbf{Q: ``There are lots of links and subpages in the webpages, what is exactly the scope
for this assignment?''}\\
A: The scope includes the links in this handout and their descendant pages that are specifically
relevant to the topics we have listed (e.g. history, events, music, food, sports). In addition, you
may also include some PDFs that can be reached from those websites, even if they are not technically
descendant pages. Use your best judgment to determine whether a webpage is relevant---a good heuristic
is whether we can ask questions about factual content included in those pages.

\vspace{0.6em}
\noindent\textbf{Q: ``Is manual scraping prohibited?''}\\
A: Manual scraping is not prohibited. To what extent you would perform the task manually is up to you.

\vspace{0.6em}
\noindent\textbf{Q: ``What libraries can I use for web scraping and document processing?''}\\
A: You can use standard libraries like Selenium, Beautiful Soup, requests, pdfminer, pypdf,
pdfplumber, and similar tools for data collection and preprocessing, as long as you provide proper
credit in your report. These are considered basic utilities rather than high-level RAG frameworks.

\vspace{0.6em}
\noindent\textbf{Q: ``What is the date range I should consider for event-based questions?''}\\
A: For any date-based questions specifically about events, we will only ask about events taking
place from October 2026 onward; every question with a specific date falls on or after October 22,
2026. Annual and recurring events are also in scope.

\vspace{0.6em}
\noindent\textbf{Q: ``Can I use a mixture-of-experts model?''}\\
A: Yes, as long as its \textbf{total} parameter count is at most 32B. We count total rather than
active parameters because you still have to load every expert, which is what limits what will run on
the hardware most of you have.

\vspace{0.6em}
\noindent\textbf{Q: ``Can I use any closed-source models (OpenAI, Claude, etc.)?''}\\
A: No. You cannot use any closed-source models for any part of the assignment, including embeddings,
retrieval, or generation. All models must be open-weight and available through HuggingFace or similar
open platforms.

\vspace{0.6em}
\noindent\textbf{Q: ``Can I use LangChain or related frameworks for this assignment?''}\\
A: No, you cannot use LangChain or other RAG frameworks (such as LlamaIndex or similar libraries) that
abstract away the core retrieval implementation. The purpose of this assignment is to learn the
underlying mechanics of retrieval systems by implementing and integrating them yourself.

\vspace{0.6em}
\noindent\textbf{Q: ``What counts as `manual implementation' for chunking and retrieval?''}\\
A: You need to write the core logic yourself. For chunking: implement the splitting logic, overlap
handling, and boundary detection. For hybrid retrieval: implement the integration pipeline, and result
processing. You can use libraries as suggested in the sections above (e.g., NLTK for tokenization,
NumPy for math operations, FAISS for vector operations, existing BM25 implementations like bm25s or
rank-bm25) but not complete retrieval frameworks like LangChain retrievers or high-level RAG pipeline
frameworks.

\vspace{0.6em}
\noindent\textbf{Q: ``How do I show that my results are statistically significant?''}\\
A: You can use statistical tests to determine whether performance differences between methods are
statistically significant. Examples include paired t-tests (for continuous metrics like F1 scores) and
McNemar's tests (for binary classification accuracy). You can run these tests on multiple metrics, or
argue in your writeup that one metric is the best proxy for final system performance. If differences
between methods are not statistically significant, interpret what this means for your system (e.g., a
new feature may not actually improve performance despite appearing better).

\vspace{0.6em}
\noindent\textbf{Q: ``If I use late days, will I be evaluated on the same test set as those who do not
use late days?''}\\
A: No. We release a different test set for each late day (up to 5), ensuring fairness for everyone
in the class. Note that the late penalty itself is the course-wide 5\% per unexcused day; see the
release policy for late days above.

%======================================================================
\problem{Acknowledgements}

\noindent This assignment was based on the Spring 2024 version of this assignment by Graham Neubig and
TAs.

%======================================================================
\problem{References}

\begin{itemize}[nosep, leftmargin=2em]
  \item Lewis et al., 2021. \href{https://arxiv.org/abs/2005.11401}{Retrieval-Augmented Generation for
    Knowledge-Intensive NLP Tasks}.
  \item Touvron et al., 2023. \href{https://arxiv.org/abs/2307.09288}{Llama 2: Open Foundation and
    Fine-Tuned Chat Models}.
  \item Vu et al., 2023. \href{https://arxiv.org/abs/2310.03214}{FreshLLMs: Refreshing Large Language
    Models with Search Engine Augmentation}.
\end{itemize}

\end{document}
